\section{Original manuscript verification and the named power operations}

\subsection*{A-029. The operators \(q_H^n\) recovered from the original scan}
The definition absent from the supplied study is explicit on the
original cote 133, scan 021, lower half (PDF page 21, printed
footer 20). Let \(G\) be a finite group, \(H\leq S_n\),
and write \(A(G)\) for its Burnside ring and
\(A_n(G)=A(G\times S_n)\). The manuscript first defines
\[
 q_H:A_n(G)\longrightarrow A(G),\qquad
                   q_H([E])=[H\backslash E],                 \tag{A-029.1}
\]
and then defines
\[
                       q_H^n(x)=q_H(\tau^n(x)).                \tag{A-029.2}
\]
The injection operation \(\tau^n\) and its completed exponential
extension are on scans 020--021 and were proved in A-011. Consequently,
for an actual finite \(G\)-set \(E\),
\[
             q_H^n([E])=[\operatorname{Inj}([n],E)/H].          \tag{A-029.3}
\]
Left versus right orbit notation is immaterial here: the \(S_n\)
action is precomposition, converted to a right action by inversion,
and commutes with the \(G\) action. These are effective operations
on actual sets and integral polynomial operations of degree \(n\)
on the additive Burnside group. They are generally neither additive
nor homogeneous under ordinary integer scaling.

\begin{proof}
Taking \(H\)-orbits commutes with disjoint union. It therefore
induces the additive homomorphism (A-029.1) on group completions.
The action of \(G\) descends because it commutes with that of \(H\).
A-011 constructs the unique homomorphism of groups
\[
 \tau^*:A(G)\longrightarrow
        1+\prod_{j\geq1}A_j(G)
\]
where multiplication is the induction product. Its degree-\(n\)
coefficient agrees with the class of injections for actual sets.
Composing that coefficient with (A-029.1) defines (A-029.2) on
all virtual classes and proves (A-029.3).

For completeness, polynomial degree is meant in ordinary additive
coordinates, with integer-valued polynomial coefficients. Choose
the finite basis \(u_1,\ldots,u_s\) of transitive \(G\)-sets
for the free abelian group \(A(G)\). For integers \(a_i\),
\[
             \tau^*\left(\sum_i a_i u_i\right)
                         =\prod_i\tau^*(u_i)^{a_i}.
\]
Write \(\tau^*(u_i)=1+v_i\), where \(v_i\) has positive
induction degree. The degree-\(n\) truncation of
\((1+v_i)^{a_i}\) is
\(\sum_{j=0}^n\binom{a_i}{j}v_i^j\), also for negative
\(a_i\), by the formal inverse identity in a complete graded ring.
A term contributing to degree \(n\) has \(\sum_i j_i\leq n\).
Its coefficient is a product of integer-valued binomial polynomials
of total degree at most \(n\). The additive map \(q_H\)
preserves this conclusion. No denominator is inverted in the
Burnside group: every binomial coefficient takes an integer value
at every integer argument.

To see that the degree is exactly \(n\), restrict to multiples
\(a\cdot1\) of the trivial point. The action of \(H\) on
injections into an \(a\)-element set is free, giving for
\(a\geq0\)
\[
 q_H^n(a\cdot1)=\frac{n!}{|H|}\binom an\,1.                    \tag{A-029.4}
\]
Both sides are polynomial functions of \(a\), so the equality
holds at every integer. The leading coefficient is \(1/|H|\),
which is nonzero over \(\mathbb Q\); thus the degree is \(n\).
For \(n=0\) the operation is the constant one, of degree zero.
For every \(n\geq2\), (A-029.4) gives \(q_H^n(1)=0\) but
\(q_H^n(n\cdot1)=n!/|H|>0\). It is therefore not homogeneous
of degree \(n\) under ordinary scaling, and is not additive.

The parenthetical wording on scan 021 appears to call the
operations polynomial and homogeneous of degree \(n\). Its
unambiguous graded meaning is that \(\tau^n\) lies in induction
degree \(n\). If ``homogeneous'' is instead interpreted as the
identity \(q_H^n(ax)=a^nq_H^n(x)\), the last calculation refutes
that interpretation. No additional homogeneous-scaling assertion
is needed for the construction.
\end{proof}

For \(H=S_n\), (A-029.3) is exactly the \(n\)-element-subset
operation; for \(H=1\), it is the ordered-injection operation.
Thus the original definition connects the supplied notation to the
subset operations already distinguished from canonical opposite
symmetric powers in A-001--A-002. Linearization gives actual
permutation representations of these orbit sets, not an automatic
exterior-power identity or a counting formula for every irreducible
Kronecker multiplicity. The previous saturation counterexamples
are unaffected.

\paragraph{Original-source verification and exact scope.}
The official source is
\url{https://grothendieck.umontpellier.fr/133.pdf}. It contains
54 PDF pages, including its cover; the catalogue's 53 is the number
of manuscript/offprint pages. In this file, the supplied scan number
is the PDF page number, while the printed footer is one smaller.
The complete PDF hash and byte count are recorded in
\path{agents/archive/GROTHENDIECK133_SOURCE_MANIFEST.json}.
In this run the following originals were visually read in full:
scans 006, 008, 009, and 019--029. Their material agrees with the
supplied transcriptions or secondary study at the specific points
used in A-001--A-015, subject to the mathematical corrections already
proved. In particular the deletions on scans 009 and 022 are visible,
the subset identities on scan 019 are not deleted, the direct-summand
claim on scan 025 really is stated for a single sufficiently large
\(H_I\), and scan 029 explicitly asks the local-ring polynomiality
question. The uncertain name on scan 008 is not authenticated.
This does not claim visual verification of the other cotes.

\paragraph{Grothendieck dictionary.}
For a finite \(G\)-set, its associated constant finite scheme carries
an action of the constant finite group scheme \(G_{\mathbb Z}\).
The injection subset and the finite orbit quotient in (A-029.3)
are actual finite schemes with the induced action. Disjoint union
supplies the additive Burnside group, and induction supplies the
graded multiplication. Linearization is the associated finite free
permutation comodule, in the representation formalism of SGA 3, I,
Proposition 4.7.2, pp.~34--35. The homogeneous induction degree
is a torus grading as in SGA 3, I, Proposition 4.7.3, pp.~35--36.
Exterior operations on locally free sheaves instead use SGA 6, V,
Definition 2.1, p.~314, and VI, Theorem 3.3, pp.~392--393.
That dictionary preserves the distinction between orbit-quotient
set operations and exterior Schur operations; it supplies no
missing specialness, positivity of virtual differences, or
Kronecker saturation theorem.
